% !PRINTED 518
\weberpage{518}

%<<CONT
gewiss eine Gruppe. Ist $\pi_\beta$ eine Permutation, durch die
die Reihe $A$ in die Reihe $B$ übergeht, die nach der
Voraussetzung der Transitivität in $P$ existiren muss, so ist
$Q_\alpha \pi_\beta$ eine Nebengruppe zu $Q_\alpha$, durch die
$A$ in $B$ übergeht. $\pi_\beta^{-1} Q_\alpha \pi_\beta =
Q_\beta$ ist eine zu $Q_\alpha$ conjugirte Gruppe, und zwar die,
durch deren Permutation die Reihe $B$ nicht verschoben wird.
Auf diese Weise erhalten wir die conjugirten Gruppen

\[ Q_\alpha,\ Q_\beta \dots Q_\sigma, \]

deren grösster gemeinschaftlicher Theiler die vorhin
betrachtete Gruppe $Q$ vom Index $j$ ist.

Eine zu $Q_\alpha$ gehörige Function $y_\alpha$ ist die Wurzel
einer Gleichung vom $s^{\mathrm{ten}}$ Grade in $\Omega$,
$\varphi(y) = 0$, und die übrigen Wurzeln $y_\beta \dots y_\sigma$
gehen aus $y_\alpha$ hervor durch Anwendung der Permutationen
der Nebengruppen $Q_\alpha \pi_\beta, \dots Q_\alpha \pi_\sigma$.
Die vorhin gefundene Gleichung $\chi(u) = 0$ vom Grade $j$ ist
die Galois'sche Resolvente der Gleichung $\varphi(y) = 0$. Wir
haben oben bemerkt, dass die Gruppe der Gleichung $\chi(u) = 0$
vom Grade $j$ isomorph ist mit der Gruppe der Permutationen,
die durch $P$ unter den Systemen der $A, B \dots S$
hervorgerufen werden. Da aber wegen der Transitivität von $P$
auch diese Systeme transitiv in einander übergehen, so muss $j$
nach dem Satze \S.\,154, 7. durch $s$ theilbar sein, und in dem
besonderen Falle, wo $j$ eine Primzahl ist, muss $j = s$ sein.

Die gleichzeitige Adjunction sämmtlicher Wurzeln von $\varphi(y)$
hat denselben Erfolg wie die Adjunction der Function $\psi$,
nämlich den, die Gruppe der Gleichung $f(x) = 0$ von $P$ auf $Q$
zu reduciren.

Durch Adjunction von $y_\alpha$ sondert sich von $f(x)$ ein
rationaler Factor $f(x, y_\alpha)$ ab, dessen Wurzeln die
Grössen der Reihe $A$ sind, und durch Adjunction der übrigen $y$
erhält man die Zerlegung

\weq{6}{f(x) = f(x,y_\alpha)\, f(x,y_\beta) \dots f(x,y_\sigma),}

und diese Functionen $f(x,y_\alpha), f(x,y_\beta) \dots
f(x,y_\sigma)$ stimmen überein mit den Factoren in (2),
$f_\alpha(x,\psi), f_\beta(x,\psi) \dots f_\sigma(x,\psi)$.

Zu bemerken ist hierbei noch, dass $f(x,y_\alpha)$ im Körper
$\Omega(y_\alpha)$ irreducibel ist. Denn wegen der
vorausgesetzten Irreducibilität von $f(x)$ giebt es in $P$
Permutationen, durch die $\alpha$ in jedes $\alpha_i$
übergeführt wird, und diese Permutationen gehören zu
$Q_\alpha$. Wenn also eine rationale Function $\varphi(x,
y_\alpha)$ für $x = \alpha$ verschwindet,
%CONT>>
